% from ega0-20-fr.tex
\oldpage[0\textsubscript{IV-1}]{123}
\begin{proposition}[20.3.6]
\label{0.20.3.6-fr}
Soient $u:A\to B$, $v:B\to C$ deux homomorphismes continus d'anneaux
topologiques, $L$ un $C$-bimodule discret annulé par un idéal bilatère ouvert
de $C$. On a une suite exacte canonique
\begin{align}
\tag{20.3.6.1}
\label{0.20.3.6.1-fr}
  0\to\operatorname{D\acute{e}r.\ cont}_B(C,L)\xrightarrow{u^0}
  \operatorname{D\acute{e}r.\ cont}_A(C,L)&\xrightarrow{v^0}
  \operatorname{D\acute{e}r.\ cont}_A(B,L)\xrightarrow{\partial} \\
  &\to\operatorname{Exantop}_B(C,L)\xrightarrow{u^1}
  \operatorname{Exantop}_A(C,L)\xrightarrow{v^1}
  \operatorname{Exantop}_A(B,L) \notag
\end{align}
où $\partial$ est défini par passage à la limite inductive à partir de
l'homomorphisme $\partial$ de \sref{0.20.2.2.1-fr} ; cette suite exacte est
fonctorielle en $L$ (dans la catégorie des $C$-bimodules discrets annulés par
des idéaux bilatères ouverts).
\end{proposition}

\begin{proof}
Cela résulte de l'exactitude du foncteur $\varinjlim$, à partir de
\sref{0.20.2.2-fr}.
\end{proof}

\begin{corollary}[20.3.7]
\label{0.20.3.7-fr}
Soient $A$, $B$, $C$ trois anneaux topologiques commutatifs, $u:A\to B$,
$v:B\to C$ deux homomorphismes continus, $L$ un $C$-module discret annulé
par un idéal ouvert de $C$. On a une suite exacte canonique de $A$-modules,
fonctorielle en $L$,
\begin{align}
\tag{20.3.7.1}
\label{0.20.3.7.1-fr}
  0\to\operatorname{D\acute{e}r.\ cont}_B(C,L)\xrightarrow{u^0}
  \operatorname{D\acute{e}r.\ cont}_A(C,L)&\xrightarrow{v^0}
  \operatorname{D\acute{e}r.\ cont}_A(B,L)\xrightarrow{\partial} \\
  &\to\operatorname{Exalcotop}_B(C,L)\xrightarrow{u^1}
  \operatorname{Exalcotop}_A(C,L)\xrightarrow{v^1}
  \operatorname{Exalcotop}_A(B,L). \notag
\end{align}
\end{corollary}

\begin{corollary}[20.3.8]
\label{0.20.3.8-fr}
Sous les hypothèses de \sref{0.20.3.5-fr} (resp.
\sref{0.20.3.6-fr}) on a une suite exacte canonique, fonctorielle en $L$
\begin{equation}
\tag{20.3.8.1}
\label{0.20.3.8.1-fr}
  0\to\operatorname{D\acute{e}r.\ cont}_B(C,L)\xrightarrow{u^0}
  \operatorname{D\acute{e}r.\ cont}_A(C,L)\xrightarrow{v^0}
  \operatorname{D\acute{e}r.\ cont}_A(B,L)\to
  \operatorname{Exantop}_{B/A}(C,L)\to0
\end{equation}
(resp.
\begin{equation}
\tag{20.3.8.2}
\label{0.20.3.8.2-fr}
  0\to\operatorname{D\acute{e}r.\ cont}_B(C,L)\xrightarrow{u^0}
  \operatorname{D\acute{e}r.\ cont}_A(C,L)\xrightarrow{v^0}
  \operatorname{D\acute{e}r.\ cont}_A(B,L)\to
  \operatorname{Exalcotop}_{B/A}(C,L)\to0\text{).}
\end{equation}
\end{corollary}

Nous laissons au lecteur le soin d'écrire les diagrammes analogues à ceux de
\sref{0.20.2.5-fr}.

\subsection{Parties principales et différentielles}
\label{subsection:0.20.4-fr}

Dans toute la suite de ce paragraphe et dans les trois suivants, tous les
anneaux sont supposés \emph{commutatifs}.

\begin{env}[20.4.1]
\label{0.20.4.1-fr}
Soient $A$ un anneau topologique, $B$ une $A$-algèbre topologique ; la
$A$-algèbre $B\otimes_A B$ sera munie de la topologie \emph{produit
tensoriel}, qui en fait une $A$-algèbre topologique ; nous désignerons par $p$
(ou $p_{B/A}$) le $A$-homomorphisme canonique surjectif
\begin{equation}
\tag{20.4.1.1}
\label{0.20.4.1.1-fr}
  p:B\otimes_A B\to B
\end{equation}
tel que $p(b\otimes b')=bb'$ ; il est immédiat que $p$ est continu. Le
\emph{noyau} de $p$ sera noté $\mathfrak{J}_{B/A}$ (ou seulement
$\mathfrak{J}$ s'il n'y a pas de confusion). On désignera par
$j_1:B\to B\otimes_A B$ et $j_2:B\to B\otimes_A B$ les deux
$A$-homomorphismes canoniques, tels que
\[
  j_1(b)=b\otimes1, \qquad j_2(b)=1\otimes b
\]
qui sont continus.
\end{env}

